\documentclass[../../main.tex]{subfiles}
\begin{document}

\begin{center}
  {\fontsize{48pt}{58pt}\selectfont \textbf{2013}\normalsize\textbf{年度}}
\end{center}

\begin{table}[htb]
  \centering
  \caption{出題テーマと難易度}
  \begin{tabular}{|c|p{9cm}|c|}
    \hline
    \textbf{問題番号} & \textbf{テーマ} & \textbf{難易度} \\
    \hline
    第1問 & 平行四辺形内の線分の比 & \\
    \hline
    第2問 & 床関数による数列の和の上界 & \\
    \hline
    第3問 & 二次式で割った余りの係数の互いに素性 & \\
    \hline
    第4問 & 余弦関数を含む式の最大値 & \\
    \hline
    第5問 & 対数曲線に接する円で囲む面積 & \\
    \hline
    第6問 & 硬貨による反射移動の確率 & \\
    \hline
  \end{tabular}
\end{table}

\begin{table}[htb]
  \centering
  \caption{解答}
  \renewcommand{\arraystretch}{1.4}
  \begin{tabular}{|c|c|p{9cm}|}
    \hline
    \textbf{問題番号} & \textbf{小問} & \textbf{答え} \\
    \hline
    第1問 & -- & $AP:PQ=19:3$ \\
    \hline
    第2問 & -- & $\sum_{n=1}^M a_n\le2^{N+1}-N-5$（証明問題） \\
    \hline
    第3問 & -- & 余り$ax+b$の係数$a,b$は互いに素な整数（証明問題） \\
    \hline
    第4問 & -- & $\dfrac{\sqrt{3}\pi^2}{16}$ \\
    \hline
    第5問 & -- & $8\sqrt{3}-2\sqrt{3}\log3-\dfrac{8\pi}{3}$ \\
    \hline
    第6問 & (1) & $\dfrac{1}{2}$ \\
    \cline{2-3}
    & (2) & $n2^{1-2n}$ \\
    \hline
  \end{tabular}
\end{table}

\paragraph{試験形態}
6問．試験時間は要確認．

\end{document}
